\documentclass[aspectratio=169]{beamer}

\usetheme{default}
\usecolortheme{default}
\setbeamertemplate{navigation symbols}{}
\setbeamertemplate{footline}[frame number]
\setbeamercolor{title}{fg=black}
\setbeamercolor{frametitle}{fg=black}
\setbeamerfont{frametitle}{series=\bfseries}

\usepackage[T1]{fontenc}
\usepackage[utf8]{inputenc}
\usepackage[brazil]{babel}
\usepackage{graphicx}
\usepackage{amsmath}
\usepackage{booktabs}
\usepackage{array}
\usepackage{tikz}
\usetikzlibrary{arrows.meta,positioning}

\title{CERISE — Dois Trabalhos, Uma Base Comum}
\subtitle{LARS 2026 (alocação de tarefas via RL) e LAFusion 2026 (fusão sensorial via EKF)}
\author{Yan Santos Leite}
\institute{Orientador: Prof. Alisson Assis Cardoso \\ Escola de Engenharia Elétrica, Mecânica e de Computação (EMC) — UFG \\ Fundação de Amparo à Pesquisa do Estado de Goiás (FAPEG)}
\date{28 de agosto de 2026}

\begin{document}

% ============================================================
\begin{frame}
\titlepage
\end{frame}

% ============================================================
\begin{frame}{Roteiro e visão geral dos dois trabalhos}
\small Dois trabalhos que compartilham a base de código e a plataforma física (três TurtleBot3 no CERISE), mas respondem perguntas diferentes. O LARS já foi submetido em 06/08; esta apresentação foca no LAFusion, a sete dias do prazo, e fecha com a conclusão e os próximos passos dos dois.

\vspace{0.1cm}
\begin{center}
\renewcommand{\arraystretch}{1.15}
\small
\begin{tabular}{p{3.0cm}p{4.4cm}p{4.4cm}}
\toprule
 & \textbf{LARS 2026} & \textbf{LAFusion 2026} \\
\midrule
Pergunta & RL supera a heurística gulosa na alocação de tarefas? & Como o EKF se comporta quando a taxa de detecção do YOLO é baixa? \\
Método & PPO (Stable-Baselines3) & Extended Kalman Filter por robô \\
Plataforma & 3 TurtleBot3, ROS2/Gazebo & 3 TurtleBot3, ROS2/Gazebo \\
Estado & Submetido no EasyChair em 06/08 & Prazo de submissão 04/09 \\
\bottomrule
\end{tabular}
\end{center}

\small A mesma câmera overhead com YOLOv8n alimenta os dois trabalhos. O LARS trata a posição detectada como dado disponível, no estado do agente. O LAFusion assume o contrário: a detecção às vezes falta, e resolver essa lacuna é o problema central.
\end{frame}

% ============================================================
\begin{frame}{Tese, argumentação e contribuições}
\small
\textbf{Tese central.} Nos dois casos, o resultado inesperado vem acompanhado da explicação de por que aconteceu.

\vspace{0.2cm}
\textbf{Como a tese se sustenta:}
\begin{itemize}
\item \textbf{LARS}: testa a hipótese de que RL supera a heurística, encontra o oposto, e refuta a explicação mais plausível a priori com um experimento causal (\emph{action masking}).
\item \textbf{LAFusion}: o EKF melhora a posição no instante da correção (+23,7\%), mas piora quando medido de forma contínua ($-$12,7\%). Isso é o comportamento esperado quando o filtro passa muito tempo sem correção, o que explica a piora.
\end{itemize}

\vspace{0.2cm}
\textbf{Contribuições:} (1) um gêmeo digital câmera+YOLO validado (mAP@0,5=0,995, erro de 4,7 cm), base comum aos dois trabalhos; (2) a causa de o PPO perder para a heurística na alocação de tarefas entre 3 robôs: o espaço de ação pequeno; (3) uma avaliação do EKF sob baixa taxa de detecção do YOLO em duas frentes: $+23,7\%$ de ganho nos instantes de correção e $-12,7\%$ de erro medido continuamente.
\end{frame}

% ============================================================
\section{LAFusion 2026}

\begin{frame}{LAFusion — Arquitetura e pipeline}
\begin{center}
\begin{tikzpicture}[
  scale=0.72, every node/.style={scale=0.72},
  node distance=4mm,
  box/.style={draw, rounded corners, align=center, minimum height=1.0cm, font=\tiny, fill=blue!6},
  arr/.style={-{Stealth[length=2mm]}, thick}
]
\node[box, minimum width=1.7cm] (cam) {Câmera\\overhead fixa};
\node[box, minimum width=2.6cm, right=of cam, fill=blue!12] (yolo) {\texttt{yolo\_detector.py}\\YOLOv8n + \texttt{projection.py}\\\texttt{/robot\_detections}};
\node[box, minimum width=2.4cm, right=of yolo] (ekf) {EKF por robô\\\texttt{ekf\_fusion\_node.py}};
\node[box, minimum width=1.7cm, right=of ekf] (pose) {Pose\\fundida};
\node[box, minimum width=1.7cm, below=9mm of ekf, fill=green!8] (odom) {Odometria\\(predição)};
\node[box, minimum width=2.6cm, above=8mm of ekf, fill=red!8] (core) {\texttt{ekf\_core.py}\\correção compartilhada\\(prod.+avaliação+sintético)};
\node[box, minimum width=2.6cm, below=9mm of core, right=8mm of core, fill=yellow!12] (assoc) {\texttt{association.py}\\gating por distância de\\Mahalanobis ($P\!+\!R$, não metros)};
\draw[arr] (cam) -- (yolo);
\draw[arr] (yolo) -- (ekf);
\draw[arr] (ekf) -- (pose);
\draw[arr] (odom) -- (ekf);
\draw[arr, dashed] (odom.east) -| (pose.south);
\draw[arr, dotted] (core) -- (ekf);
\draw[arr, dotted] (assoc) -- (ekf);
\end{tikzpicture}
\end{center}

\vspace{0.1cm}
\scriptsize\texttt{yolo\_detector.py} já converte pixel$\to$mundo (via \texttt{projection.py}) antes de publicar \texttt{/robot\_detections}. A seta tracejada mostra que o EKF publica pose a cada leitura de odometria (predição), não só após uma correção YOLO. \texttt{association.py} e \texttt{ekf\_core.py} não são nós ROS2 à parte: são módulos importados e chamados de dentro do próprio \texttt{ekf\_fusion\_node.py} (setas pontilhadas). \texttt{association.py} faz o gating de Mahalanobis por vizinho mais próximo (não JPDA nem Hungarian), suficiente para $N=3$ mas com risco de troca de identidade se dois robôs ficarem muito próximos. \texttt{ekf\_core.py} concentra a correção (forma de Joseph), reaproveitada por produção, avaliação offline e validação sintética.

\vspace{0.1cm}
\small\textbf{Por que EKF, e não UKF ou filtro de partículas:} o modelo de observação da câmera é quase linear (projeção pinhole numa cena planar), então o ganho de um UKF seria marginal. O problema central aqui é a \emph{taxa} de observação, não a forma do modelo.
\end{frame}

% ============================================================
\begin{frame}{LAFusion — Design do EKF e validação}
\begin{columns}[T]
\begin{column}{0.54\textwidth}
\scriptsize
\textbf{Estado:} $x_k=[p_x,p_y,\theta]^\top$ por robô. A predição integra o \emph{delta} de odometria, não o valor absoluto, preservando correções visuais anteriores.

\vspace{0.1cm}
\textbf{Hiperparâmetros e sua origem:}
\begin{itemize}\scriptsize\itemsep1pt
\item $Q=\text{diag}(0{,}5,\,0{,}5,\,0{,}25)$: calibrado contra dados reais, ganho satura em $Q\geq0{,}5$.
\item Teto $\text{clip}(P,-\infty,0{,}05)$: sem ele $P$ cresce sem limite (traço $\approx$85) e quebra o gating, deixando um robô ``roubar'' a detecção do vizinho.
\item Gating $d^2\leq9{,}21$: padrão do $\chi^2$ com 2 g.l.\ a 99\%.
\item Forma de Joseph e gating com $S=P+R$: rigor numérico, reproduzem os resultados publicados.
\end{itemize}

\vspace{0.1cm}
\emph{Por que o teto existe:} mantém a associação estável, mas o estado segue impreciso sem correção. Um \emph{fading factor} (AFKF), teoricamente melhor, foi testado e perde: o ganho sob drift agressivo cai de $+23{,}7\%$ para $-16{,}5\%$.
\end{column}
\begin{column}{0.46\textwidth}
\scriptsize
\textbf{Validação sintética, antes de qualquer dado real:} Monte Carlo, 30 sementes, 2000 passos, contra verdade terrestre conhecida.

\vspace{0.1cm}
\begin{center}
\begin{tabular}{lcc}
\toprule
\textbf{Estat.} & \textbf{Obtido} & \textbf{Esperado} \\
\midrule
NEES & 2,955 & 3,0 (1,5\%) \\
NIS & 1,982 & 2,0 (0,9\%) \\
\bottomrule
\end{tabular}
\end{center}

\vspace{0.1cm}
Isso isola bugs de implementação de problemas de dado: um filtro inconsistente contra dados sintéticos jamais seria confiável contra dados reais.

\vspace{0.1cm}
\textbf{Dados reais:} 3 cenários no Gazebo (parado, reto, curva/oclusão), 3 robôs, sincronização verificada.

\vspace{0.1cm}
\textbf{Regressão produção vs.\ avaliação:} \texttt{predict()} não é compartilhada entre os dois (modelos de movimento diferentes), então um teste compara numericamente as duas implementações a cada passo, para os mesmos dados de entrada, e falha se divergirem.
\end{column}
\end{columns}
\end{frame}

% ============================================================
\begin{frame}{O mesmo filtro ganha $+23{,}7\%$ na correção e perde $-12{,}7\%$ medido continuamente}
\begin{columns}[T]
\begin{column}{0.5\textwidth}
\centering\scriptsize
\textbf{(a) Nos instantes de correção YOLO}\\[2pt]
\begin{tabular}{lccc}
\toprule
\textbf{Drift} & \textbf{EKF/Odom} & \textbf{$\Delta$} & \textbf{$\Delta$RMSE} \\
\midrule
Nenhum & 0,036/0,000 & N/A & N/A \\
Leve & 0,053/0,050 & $-5,3\%$ & $-2,1\%$ \\
Agressivo & 0,128/0,168 & \textbf{+23,7\%} & \textbf{+16,7\%} \\
\bottomrule
\end{tabular}\\[2pt]
\raggedright $n=1323$, Wilcoxon $p<0{,}001$ sob drift agressivo.

\vspace{0.1cm}
\emph{Por que o ganho é negativo sob drift leve:} o erro de detecção visual ($\approx$3–6 cm) ainda supera o de uma odometria pouco corrompida, então fundir as duas fontes piora o resultado. É o comportamento esperado de um filtro bem calibrado.
\end{column}
\begin{column}{0.5\textwidth}
\centering\scriptsize
\textbf{(b) A cada leitura de odometria}\\[2pt]
\begin{tabular}{lccc}
\toprule
\textbf{Cenário} & \textbf{Cobert.} & \textbf{EKF/Odom} & \textbf{$\Delta$} \\
\midrule
Parado & 27,6\% & 0,212/0,177 & $-19,3\%$ \\
Reto & 11,3\% & 0,212/0,178 & $-18,8\%$ \\
Curva/ocl. & 0,0\% & 0,178/0,178 & $-0,0\%$ \\
\textbf{Agregado} & --- & \textbf{0,200/0,178} & \textbf{$-12,7\%$} \\
\bottomrule
\end{tabular}\\[2pt]
\raggedright $n=1607$, Wilcoxon $p<0{,}001$.

\vspace{0.1cm}
Com 0\% de cobertura a mudança é exatamente 0,0\%: sem correções o EKF degenera para integração pura de delta-odometria. A coincidência exata confirma o mecanismo.
\end{column}
\end{columns}

\vspace{0.1cm}
\scriptsize\textbf{ATE vs.\ RPE (Sturm et al.\ 2012):} o erro em (b) é ATE, absoluto por amostra; o RPE (drift entre passos consecutivos) piora bem menos ($-4,2\%$), então a reversão vem de erro \emph{acumulando} sob correção esparsa. \textbf{Fundamentação:} Sinopoli et al.\ (2004) mostram que a covariância diverge abaixo de uma taxa crítica de observação, e 0\% de cobertura é o caso-limite $p\to0$; o teto evita o crescimento ilimitado, mas limita só a \emph{incerteza estimada}, deixando o filtro confiante demais na predição só-odometria. Dissanayake et al.\ (2001) provam convergência do EKF-SLAM sob observação \emph{densa}, premissa violada aqui por construção. UKF, partículas e IMM não resolvem: a causa raiz é a topologia do sensor, não a escolha do filtro.
\end{frame}

% ============================================================
\begin{frame}{Referências-chave e próximos passos}
\begin{columns}[T]
\begin{column}{0.5\textwidth}
\textbf{LARS — bibliografia essencial:}
\tiny
\begin{itemize}
\item Henderson et al. (2018): reprodutibilidade em RL profundo, motiva o teste pareado.
\item Huang \& Ontañón (2022): action masking produz gradiente válido; no nosso domínio elimina a violação sem fechar o gap.
\item Agrawal et al. (2023, RTAW), Paul et al. (2022, Capsule Attention): arquiteturas atencionais como direção futura.
\end{itemize}
\vspace{0.1cm}
\textbf{Próximos passos:} aguardar aceite (15/09/2026); espaço de ação e orçamento de treino maiores; formulação Dec-POMDP (CTDE).
\end{column}
\begin{column}{0.5\textwidth}
\textbf{LAFusion — bibliografia essencial:}
\tiny
\begin{itemize}
\item Sinopoli (2004), Bailey \& Durrant-Whyte (2006), Dissanayake et al. (2001): fundamentação teórica clássica do achado central (erro contínuo, modelo EKF, convergência sob observação densa).
\item Khosravi et al. (2026): localização cooperativa com fusão assíncrona, cenário mais próximo do nosso.
\item Ahmed \& Frémont (2026): rastreio multi-robô descentralizado, mesma classe de problema.
\item Dobrynin \& Garcia (2025): testbed TurtleBot3 com EKF distribuído.
\end{itemize}
\vspace{0.1cm}
\textbf{Próximos passos:} submissão no CMT até 04/09/2026; $Q$ adaptativo por tempo desde a última correção; Jacobiano completo na projeção; associação robusta (Hungarian/JPDA); LiDAR como terceira fonte.
\end{column}
\end{columns}

\vspace{0.1cm}
\hrule
\vspace{0.1cm}
\tiny
\textbf{Estado do LAFusion:} 10 páginas escritas (limite do CFP é 12), prazo \textbf{04/09/2026}, Springer CCIS, double-blind via CMT. Em 16/08 a calibração da câmera foi corrigida: com os 4 intrínsecos livres estava mal-condicionada ($f_x$ 33\% acima do esperado), o caso degenerado de Zhang (2000), já que a câmera fixa deixa o tabuleiro sempre fronto-paralelo. Fixamos os intrínsecos no valor nominal e calibramos só a distorção; o resultado reproduz a projeção de produção a menos de 0,001\,m e \textbf{nenhum resultado publicado muda}.

\vspace{0.05cm}
\textbf{Conclusão:} no LARS, o \emph{action masking} refuta a explicação mais plausível a priori para o resultado negativo e redireciona o diagnóstico para o espaço de ação e o orçamento de treino. No LAFusion, o ganho de +23,7\% na correção convive com uma perda de $-$12,7\% medida continuamente, e a teoria de filtragem sob observação intermitente explica a reversão. Os dois trabalhos reportam o resultado completo e o fundamentam.
\end{frame}

\end{document}
